% Copyright (c) 2026 ArkLib Contributors. All rights reserved.
% Released under Apache 2.0 license as described in the file LICENSE.

\section{The Plonk Polynomial IOP}\label{sec:plonk}

ArkLib formalizes the polynomial oracle reduction underlying Plonk~\cite{GWZC19}. It
composes four stages, each a standalone oracle reduction. The result concerns the
polynomial IOP; it does not include a polynomial commitment scheme or Fiat--Shamir.

Fix a field $\field$ with decidable equality, a number of rows $n\geq 2$, and a primitive
$n$-th root of unity $\omega$, so that $H=\{1,\omega,\ldots,\omega^{n-1}\}$ is the row
domain. A constraint system has $m$ variables and $n$ gates; each gate constrains a left,
right and output wire, and copy constraints identify wires that carry the same variable.
Fix coset labels for the three wire columns and a public-input polynomial $\mathit{pi}$ of
degree less than $n$. A wire table is three polynomials $a_0,a_1,a_2$ of degree less than
$n$.

\begin{definition}[Plonk oracle reduction]\label{def:plonk_assembly}
\lean{Plonk.Assembly.oracleReduction, Plonk.Assembly.rel,
Plonk.GrandProduct.InputRelation}
\leanok
\uses{def:sendsinglewitness_oracle_reduction,def:generator_batch,def:quotient_zero_test}
The input relation holds for a wire table when the gate polynomial built from the table and
$\mathit{pi}$ vanishes on $H$, and the wire values satisfy every copy constraint. The
reduction is the sequential composition of four stages, with six rounds in total.
\begin{enumerate}
  \item \textbf{Wire reveal.} The prover sends the wire table as one polynomial oracle.
  \item \textbf{Grand product.} The verifier sends the permutation challenges
    $(\beta,\gamma)$ jointly. The prover replies with the accumulator polynomial $Z$.
    The output claim is that three residual polynomials vanish on $H$.
  \item \textbf{Power batch.} The verifier sends one challenge $\alpha$ and the three
    residuals are combined with coefficients $1,\alpha,\alpha^2$
    (Definition~\ref{def:generator_batch}).
  \item \textbf{Quotient.} The prover sends the quotient polynomial and the verifier
    checks the quotient identity at a random point
    (Definition~\ref{def:quotient_zero_test}). The output relation accepts every output.
\end{enumerate}
Every earlier oracle is retained up to the terminal check.
\end{definition}

\begin{theorem}[Plonk completeness]\label{thm:plonk_completeness}
\lean{Plonk.Assembly.oracleReduction_completeness}
\leanok
\uses{def:plonk_assembly}
Let $\field$ be finite and sampleable. For every initial distribution of the shared oracle
state and every implementation of that oracle, the honest prover is complete from the input
relation to the output relation with error $3n/|\field|$. The failure occurs only at the
grand-product challenge.
\end{theorem}

\begin{theorem}[Plonk knowledge soundness]\label{thm:plonk_knowledge_soundness}
\lean{Plonk.Assembly.oracleVerifier_rbrKnowledgeSoundnessWith, Plonk.Assembly.error,
Plonk.Assembly.sum_error, Plonk.Assembly.transcriptExtractor,
Plonk.Assembly.oracleVerifier_knowledgeSoundnessWith}
\leanok
\uses{def:plonk_assembly}
Under the hypotheses of Theorem~\ref{thm:plonk_completeness}, the verifier is
round-by-round knowledge sound from the input relation to the output relation. The errors
at the three challenges are
\[
  \frac{3n-1}{|\field|}\ \text{(grand product)},\qquad
  \frac{2}{|\field|}\ \text{(power batch)},\qquad
  \frac{4(n-1)}{|\field|}\ \text{(quotient)},
\]
and the reveal stage has no challenge. The errors sum to $(7n-3)/|\field|$. Consequently
the verifier is knowledge sound with error $(7n-3)/|\field|$, and the extractor returns
the wire table from the first prover message.
\end{theorem}

\paragraph{The circuit relation.}
For a constraint system with $\ell\leq m$ public inputs, the Plonk relation holds for a
public input $x\in\field^{\ell}$ and a private witness $w\in\field^{m-\ell}$ when the
assignment $(x,w)$ satisfies every gate and copy constraint. To prove it, the constraint
system is prepared: $\ell$ public-equality rows are prepended, so the domain has $\ell+n$
rows and $\omega$ is a primitive $(\ell+n)$-th root of unity with $\ell+n\geq 2$. The
public constants of the new rows are cleared and $x$ enters through the public-input
polynomial.

\begin{theorem}[Plonk knowledge soundness for the circuit relation]
\label{thm:plonk_circuit_knowledge_soundness}
\lean{Plonk.relation, Plonk.Assembly.Circuit.oracleReduction, Plonk.Assembly.Circuit.relIn,
Plonk.Assembly.Circuit.oracleVerifier_rbrKnowledgeSoundnessWith,
Plonk.Assembly.Circuit.oracleVerifier_knowledgeSoundnessWith}
\leanok
\uses{thm:plonk_knowledge_soundness}
With the prepared system above, the verifier is knowledge sound from the Plonk relation at
public input $x$ to the accepting output relation, with error
$(7(\ell+n)-3)/|\field|$. The extractor decodes the first prover message into the private
witness $w$.
\end{theorem}
